%%=======================================================================
%% paper.tex
%% Title : From Natural Language to Verifiable Reasoning:
%%         A Semantic Reduction Framework for Reliable AI Mathematics
%%
%% C2 challenge deliverable. Structure follows the mandated outline
%% (materials/中文论文大纲（AI4Math）.pdf), sections 1-9.
%%
%% Compiles with: pdflatex -> bibtex -> pdflatex -> pdflatex
%%=======================================================================
\documentclass[11pt,a4paper]{article}

%% ---- packages ----------------------------------------------------------
\usepackage[T1]{fontenc}
\usepackage[utf8]{inputenc}
\usepackage{lmodern}
\usepackage{amsmath,amssymb,amsthm,mathtools}
\usepackage{stmaryrd}   % provides \llbracket / \rrbracket used in Def. 1
\usepackage{graphicx}
\usepackage{booktabs}
\usepackage{array}
\usepackage{tabularx}
\usepackage[margin=2.4cm]{geometry}
\usepackage[hidelinks]{hyperref}
\usepackage{xcolor}
\usepackage{caption}
\usepackage{enumitem}
\usepackage{microtype}
\usepackage{natbib}

%% ---- theorem environments ---------------------------------------------
\theoremstyle{plain}
\newtheorem{definition}{Definition}
\newtheorem{proposition}{Proposition}
\theoremstyle{definition}
\newtheorem{example}{Example}

%% ---- macros ------------------------------------------------------------
%% \RIS is used both in text and inside table headers (which are already in a
%% bold/roman context). \textnormal guards against nested font-shape requests
%% such as "T1/lmr/m/scit undefined" when \RIS lands inside italics.
\newcommand{\RIS}{\textnormal{\textsc{RIS}}}
\newcommand{\OK}{\ensuremath{\mathsf{ok}}}      % math-safe: used inside $...$
\newcommand{\FAIL}{\ensuremath{\mathsf{fail}}}
\newcommand{\LayerG}{\mathcal{L}_{\mathrm{gen}}}
\newcommand{\LayerR}{\mathcal{L}_{\mathrm{red}}}
\newcommand{\LayerV}{\mathcal{L}_{\mathrm{ver}}}
\newcommand{\code}[1]{\texttt{#1}}

%% ---- tight lists -------------------------------------------------------
\setlist{nosep,leftmargin=*}

\begin{document}

%%=======================================================================
\title{\vspace{-1.2cm}\bfseries From Natural Language to Verifiable Reasoning:\\
A Semantic Reduction Framework for Reliable AI Mathematics}
\author{}\date{}
%%=======================================================================

\begin{abstract}
Large language models produce mathematical prose that is fluent, plausible
and frequently wrong. The dominant response in the literature is
\emph{post-hoc verification}: generate a candidate proof, then attempt to
check it. We argue that this response is structurally incomplete, because it
treats verification as a function of the generator's output alone. We
introduce a three-layer model of neural mathematical reasoning---generation,
\emph{semantic reduction}, and verification---and identify the reduction
layer as the binding constraint on reliability.

Our contribution is threefold. First, we give a formal account of semantic
reduction as a partial, many-to-one map between natural-language
justifications and typed intermediate representations, and we show that
the composition ``generate $\rightarrow$ reduce $\rightarrow$ verify'' is
\emph{not} monotone in the quality of the generator: improving Layer 1 can
leave aggregate reliability unchanged or even degrade it, because a more
fluent generator produces justifications whose intended reading is more
difficult to recover. Second, we introduce a four-class taxonomy of
reduction defects (undertyping, scope drift, unbound symbol, non-local
warrant) and a mechanically checkable metric, the Reduction Integrity Score
(\RIS{}), that is computable from a formal artifact alone. Third, we report
a reproducible probe study ($n=108$) in which \RIS{} separates clean from
defective reductions and tracks problem difficulty monotonically. We are
explicit about the study's limits: the corpus is \emph{constructed} to
validate the instrument, and our numbers are not field estimates of error
rates. Our thesis is that reliability engineering in AI4Math should target
the reduction layer directly, and we give a concrete reliability checklist
that follows from the analysis.
\end{abstract}

%%=======================================================================
\section{Introduction}
\label{sec:intro}

Large language models have become fluent producers of mathematical
argument \citep{brown2020}. Chain-of-thought prompting made intermediate
reasoning steps explicit and improved multi-step accuracy substantially
\citep{wei2022cot}; program-of-thought prompting moved computation out of
the token stream entirely, delegating arithmetic to an interpreter
\citep{chen2022pot}; and tool-augmented models learned to call external
APIs \citep{schick2023toolformer}. These results share a premise: if a
model is allowed to externalise its reasoning into a form a machine can
execute, that reasoning can be checked \citep{alphageometry2024}.

The premise is only half right, and the half that fails is where the
reliability problem lives. A natural-language justification---the form
almost all of these systems actually produce---is not directly verifiable.
Nothing in the token stream carries types, binds symbols, or fixes the scope
of a quantifier. Verification therefore requires an intermediate step that
converts informal text into a formal object, and that step is a translation
whose failure modes are not visible in the text it consumes.

We make this precise. Let $\mathcal{N}$ be the space of natural-language
justifications, $\mathcal{F}$ the space of well-typed formal artifacts, and
let $V\colon\mathcal{F}\to\{\OK,\FAIL\}$ be a verifier such
as a proof assistant's kernel \citep{lean4tool,isabelle2002,coqmanual}. A
system's end-to-end reliability is not $\mathrm{Acc}(\LayerG)$ nor
$\Pr[V=\OK]$, but
\begin{equation}
\label{eq:reliability}
R \;=\; \Pr\!\big[\,V(\rho(n))=\OK\ \wedge\ \mathrm{Valid}(\rho(n))\,\big],
\qquad
n\sim\LayerG,\quad \rho\colon\mathcal{N}\rightharpoonup\mathcal{F},
\end{equation}
where $\rho$ is the reduction map, $\rightharpoonup$ indicating that $\rho$
is \emph{partial} and many-to-one, and $\mathrm{Valid}$ denotes agreement
between the artifact and the proof actually intended. Equation
\eqref{eq:reliability} is the paper's central formal object: reliability is
a property of the \emph{composition}, and $\rho$ appears inside the
probability, not outside it.

The practical consequence is a claim we call the \textbf{reduction
bottleneck hypothesis}: in current systems, aggregate reliability is
constrained primarily by the quality of $\rho$, not by the fluency of
$\LayerG$ nor by the strictness of $V$. Section~\ref{sec:bottleneck}
develops this; Section~\ref{sec:framework} gives the machinery.

\paragraph{Contributions.}
\begin{enumerate}[label=(\roman*)]
\item A three-layer formalisation of neural mathematical reasoning in which
      semantic reduction is a first-class, typed component rather than an
      implementation detail (Section~\ref{sec:model}).
\item A proof that verifier strength alone is an unreliable lever, because
      a perfect verifier over a lossy reduction certifies a distorted claim
      (Proposition~\ref{prop:nonmonotone}).
\item A four-class defect taxonomy for reduction---undertyping, scope drift,
      unbound symbol, non-local warrant---together with \RIS{}, a metric
      computable from the artifact alone (Section~\ref{sec:ir}).
\item A reproducible probe study over $108$ constructed items, with all
      numbers emitted by committed code, plus an explicit statement of what
      the study does \emph{not} establish (Section~\ref{sec:exp}).
\item A reliability checklist derived from the analysis
      (Section~\ref{sec:discussion}).
\end{enumerate}

We emphasise at the outset what makes this a reliability paper rather than
a mathematics paper. We do not prove new theorems about $\mathbb{N}$. We
argue that the binding constraint in AI4Math is a \emph{systems} problem
sitting between two components that the field already has, and we make that
argument falsifiable by supplying a metric and a measurement procedure.

%%=======================================================================
\section{A Three-Layer Decomposition}
\label{sec:model}

\subsection{Generation layer}

$\LayerG$ maps a problem statement $p$ to a distribution over
justifications, $n\sim\LayerG(\cdot\mid p)$. Its characteristic properties
are probabilistic, untyped and informal. The untypedness is not an oversight
but a consequence of the training objective: next-token likelihood rewards
locally plausible continuations, not well-formed bindings
\citep{bender2021parrots}. A representative artefact is a chain of
thought \citep{wei2022cot}.

\subsection{Semantic reduction layer}

$\rho\colon\mathcal{N}\rightharpoonup\mathcal{F}$ maps a justification to a
formal artifact. This is the layer this paper is about, and it deserves a
precise characterisation.

\begin{definition}[Typed intermediate representation]
\label{def:tir}
A \emph{typed intermediate representation} (TIR) is a triple
$\Phi=\langle\mathcal{T},\llbracket\cdot\rrbracket,\mathcal{C}\rangle$
where $\mathcal{T}$ is a set of \emph{types}, $\llbracket\cdot\rrbracket$
is an \emph{operational semantics} assigning to each $\varphi\in\Phi$ a
denotation $\llbracket\varphi\rrbracket$, and $\mathcal{C}$ is a set of
\emph{side conditions} that must hold for the artifact to elaborate.
Elaboration succeeds iff $\varphi$ is well-typed under $\mathcal{T}$ and
satisfies $\mathcal{C}$.
\end{definition}

Definition~\ref{def:tir} is the Curry--Howard idea used as a
\emph{engineering contract}: the artifact is not a string, it is a
term in a typed language whose typing judgement carries the information a
verifier needs. Lean \citep{lean4tool}, Isabelle/HOL \citep{isabelle2002}
and Coq \citep{coqmanual} instantiate $\mathcal{T}$, $\llbracket\cdot\rrbracket$
and $\mathcal{C}$ concretely.

$\rho$ has three properties that jointly explain why it is the bottleneck.
It is \textbf{partial} (some justifications admit no well-typed artifact);
\textbf{many-to-one} (distinct justifications may reduce to the same
artifact, so reduction discards information); and \textbf{not
syntax-preserving} (the natural-language reading and the artifact reading
can diverge without any local signal).

\subsection{Verification layer}

$\LayerV$ decides whether an artifact establishes its goal. We follow the
outline's four checks: self-consistency, derivability in a theory,
consistency with stated facts, and satisfaction of constraints. The
guarantee $\LayerV$ offers is total: it never accepts a malformed artifact.
This strength is exactly why its weakness is easy to overlook---a verifier
can only certify what it is handed.

%%=======================================================================
\section{The Bottleneck: Why Reduction, Not Generation}
\label{sec:bottleneck}

\subsection{Verification strength is not a monotone lever}

\begin{proposition}[Non-monotonicity of end-to-end reliability]
\label{prop:nonmonotone}
Fix $n$ and let $\rho_1,\rho_2$ be two reduction maps with
$\mathrm{Valid}(\rho_1(n))\geq\mathrm{Valid}(\rho_2(n))$. There exists a
verifier $V$ and a justification $n$ with $V(\rho_1(n))=\OK$ and
$V(\rho_2(n))=\OK$, yet $\mathrm{Valid}(\rho_2(n))=0$.
\end{proposition}

\begin{proof}[Sketch]
Let $n$ state a true claim $C$ with a justification whose explicit
statement is of a different, weaker claim $C'$ with $C\Rightarrow C'$. Take
$\rho_2$ to select the literal reading of $n$, yielding an artifact proving
$C'$. Then $\mathrm{Valid}(\rho_2(n))=0$ because the artifact does not
establish the intended claim $C$, while any type-correct $V$ accepts
$\rho_2(n)$ as a well-formed proof of $C'$. Hence strengthening $V$ cannot
repair a reduction error, because the defect is already baked into the
artifact before $V$ is consulted. A fortiori, raising the quality of
$\LayerG$ does not help when the error lies in $\rho$.
\end{proof}

Proposition~\ref{prop:nonmonotone} is deliberately mild, and that is the
point: it is an elementary observation, not a deep theorem. Its rhetorical
force is that it identifies a failure mode which \emph{no} amount of
downstream checking can reach, and which therefore has to be fixed at the
reduction layer. The history of computer-assisted proofs offers a parallel:
such proofs became widely accepted only when the computation itself became
independently checkable, rather than merely reported as correct. The same
structural move---from assertion to checkable artifact---is what we are
proposing for the reduction step.

\subsection{Where the field has invested, and what that implies}

The dominant paradigm couples a generator to a formal prover and reports
end-to-end success \citep{alphaproof2025,alphageometry2024,dsp2022}. This
is a major advance, and we do not dispute it. But note where the
architectural weight falls: the artifact is produced \emph{already formal}
because autoformalization is folded into the model, so $\rho$ is not
measured, reported, or debugged. We argue this hides the very quantity that
predicts downstream success. Our own evidence for this is limited---we
demonstrate the mechanism, we do not measure it in a production system---and
Section~\ref{sec:limits} states the gap plainly.

\subsection{A measurable consequence}

If reduction is the constraint, then a pipeline should exhibit a
characteristic signature: a substantial fraction of generations that are
\emph{plausible but unreducible}. Such items are invisible to end-to-end
accuracy, because the pipeline reports failure and the operator cannot
distinguish ``the model could not solve this'' from ``the model solved it in
a language nobody can type-check.'' We quantify the size of this
uninformative-failure region in Section~\ref{sec:exp} via $\RIS$.

%%=======================================================================
\section{The Proposed Framework}
\label{sec:framework}

\subsection{The reduction pipeline}

\begin{figure}[htbp]
\centering
\fbox{\begin{minipage}{0.92\textwidth}
\centering
\textbf{Layer 1 --- Generation}\\[2pt]
$\mathcal{N}$: natural-language justification\\[6pt]
$\Downarrow\quad \rho\ \text{(semantic reduction)}$\\[6pt]
\textbf{Layer 2 --- Typed IR $\Phi=\langle\mathcal{T},\llbracket\cdot\rrbracket,\mathcal{C}\rangle$}\\[6pt]
$\Downarrow\quad V\colon\mathcal{F}\to\{\OK,\FAIL\}$\\
\textbf{Layer 3 --- Verification}\\[6pt]
$\Downarrow$\\
$\mathcal{N}$: rendered explanation $\rho^{\dagger}$
\end{minipage}}
\caption{The three-layer pipeline. Natural language is treated as an
\emph{interface} to a typed core, not as the object of verification. The
renderer $\rho^{\dagger}$ is needed for interpretability and is itself a
partial inverse of $\rho$; we discuss its (non-)existence in
Section~\ref{sec:bidir}.}
\label{fig:pipeline}
\end{figure}

Figure~\ref{fig:pipeline} shows the architectural claim. The three
properties of $\rho$ identified above are addressed by three design
requirements: because $\rho$ is partial, its failure must be
\emph{reportable} rather than silent; because it is many-to-one, the
artifact must carry \emph{evidence} linking it back to the source text; and
because it is not syntax-preserving, the round trip must be
\emph{checkable}. The framework's contribution is a concrete instrument for
each.

\subsection{Reporting reduction failure: the \RIS{} metric}
\label{sec:ir}

The central instrument is a scalar computable from the artifact.

\begin{definition}[Reduction defect]
\label{def:defect}
A \emph{reduction defect} is a condition under which $\rho(n)=\varphi$ is
well-formed in isolation yet fails to represent $n$. We distinguish four
classes, one per failure mode of Definition~\ref{def:tir}:
\begin{description}[leftmargin=2.2em,labelsep=0.5em]
\item[A.] \textbf{Undertyping} --- $\varphi$ is used at a type not declared
      in $\mathcal{T}$, so the denotation $\llbracket\varphi\rrbracket$ is
      undefined.
\item[B.] \textbf{Scope drift} --- the quantifier structure of $n$ is
      narrowed or widened in $\varphi$, changing which instances the claim
      ranges over.
\item[C.] \textbf{Unbound symbol} --- $\varphi$ references an identifier
      outside its binding context.
\item[D.] \textbf{Non-local warrant} --- a step in $\varphi$ is justified
      only by material outside $\Phi$, so $\mathcal{C}$ cannot be checked
      locally.
\end{description}
\end{definition}

\begin{definition}[Reduction Integrity Score]
\label{def:ris}
Let $U$ be the set of symbols used in $\varphi$ and $B$ the set bound in
$\varphi$. Define coverage $\kappa = 1 - \lvert U\setminus B\rvert/\lvert U\rvert$
(with $\kappa=1$ when $U=\varnothing$) and let
$\delta(\varphi)\in\{0,\dots,4\}$ be the number of defect classes realised.
The \emph{Reduction Integrity Score} is
\begin{equation}
\label{eq:ris}
\RIS(\varphi) \;=\; 0.40\,\mathbb{1}[\varphi \text{ elaborates}]
\;+\; 0.35\,\kappa(\varphi)
\;+\; 0.25\left(1 - \frac{\delta(\varphi)}{4}\right).
\end{equation}
\end{definition}

Two properties make \RIS{} usable. First, it is \emph{mechanically
computable}: elaboration is decided by the prover's kernel, and $U$, $B$ are
syntactic. Second, it is \emph{monotone decreasing} in each defect count, so
a pipeline can use it for triage. The weights in
Equation~\eqref{eq:ris} were fixed \emph{a priori} and deliberately not
tuned on the evaluation data; tuning them on the evaluation set would make
the resulting numbers uninterpretable, so we do not do it. A reader who
disagrees with the weights can rerun the analysis with different ones, and
the code supports that.

\begin{proposition}[Bounds]
\label{prop:ris}
$\RIS(\varphi)\in[0,1]$ for all $\varphi$, with $\RIS(\varphi)=1$ iff
$\varphi$ elaborates, binds every symbol it uses, and realises no defect.
\end{proposition}

\begin{proof}
Each of the three terms is in $[0,1]$ and the weights are non-negative and
sum to one. Equality holds iff all three terms equal one, i.e.\ elaboration
succeeds ($\kappa=1$), and $\delta=0$.
\end{proof}

\subsection{Bidirectional alignment}
\label{sec:bidir}

Since $\rho$ is many-to-one, a strict inverse $\rho^{-1}$ does not exist:
the renderer cannot recover which of several source readings was intended.
We therefore do not claim one. Instead we require a
\emph{round-trip witness}: a rendering $\rho^{\dagger}(\varphi)$ together
with a check that re-reducing it returns $\varphi$. This is weaker than
inversion---it certifies consistency of the pair rather than recovery of
the original intent---and it is checkable, which inversion is not. The
residual gap between intended and rendered meaning is precisely what
Proposition~\ref{prop:nonmonotone} exploits, and we name it the
\emph{semantic residual} and leave quantifying it as future work.

%%=======================================================================
\section{Related Work and Positioning}
\label{sec:related}

\paragraph{Neural reasoning.} Chain-of-thought \citep{wei2022cot},
program-of-thought \citep{chen2022pot}, tool use \citep{schick2023toolformer}
and few-shot prompting \citep{brown2020} improve reasoning by changing what
the model emits. They are orthogonal to $\rho$: each can be paired with
whatever reduction map the downstream prover requires.

\paragraph{Neuro-symbolic and formally verified reasoning.} AlphaGeometry
\citep{alphageometry2024} couples a language model to a symbolic deduction
engine; AlphaProof \citep{alphaproof2025} pairs autoformalization with
reinforcement learning inside Lean; Draft--Sketch--Prove
\citep{dsp2022} uses informal proofs to steer a formal prover; and
autoformalization corpora such as Lean Workbook \citep{leanworkbook2024}
supply training signal. Our position is a \emph{diagnostic} one: we agree
that formal anchors are necessary, and argue they are currently
unmeasured. $\rho$ is implicit in this line of work, and making it explicit
is our contribution.

\paragraph{Evaluation.} MiniF2F \citep{minif2f2021} measures
cross-system formal reasoning; MATH \citep{hendrycks2021math} and
GSM8K \citep{cobbe2021verifiers} measure informal mathematical ability, the
latter by training verifiers directly. \RIS{} complements these by scoring
the \emph{transition} rather than either endpoint.

\paragraph{Reliability.} Work on model limitations \citep{bender2021parrots}
and on representation structure \citep{huh2024platonic} frames why fluent
output is not trustworthy output. The Platonic representation hypothesis
in particular suggests that failures concentrate in mismatches between
representational formats---which is, structurally, the same claim we make
about $\rho$, stated in representational rather than type-theoretic terms.

\paragraph{Mechanism studies.} Work analysing what models actually compute
internally \citep{kaprison2025physics} provides a complementary view:
where we reason forward from interface design, such analyses reason backward
from activations. We regard the two as compatible and jointly necessary.

%%=======================================================================
\section{Applications}
\label{sec:apps}

The framework is prescriptive in three settings.

\paragraph{Automated theorem proving.} Generate a proof sketch, reduce it to
a TIR, and require elaboration before search begins. The practical benefit
is diagnostic: when search fails, $\RIS{}$ distinguishes ``no valid artifact
was produced'' from ``valid artifact, no proof found.'' These demand
completely different responses, and current pipelines conflate them.

\paragraph{Mathematical problem solving.} For word problems, reduce to a
symbolic program and delegate evaluation to a solver \citep{chen2022pot}.
Here $\RIS$ acts as an early filter: an item whose program does not bind its
symbols should be rejected before the solver's answer is trusted.

\paragraph{Assessment.} Student solutions reduce to the same TIR, so grading
becomes a type-elaboration question rather than an assessment of fluency
\citep{schick2023toolformer}. We regard this as the highest-value
application and flag its ethical requirements in Section~\ref{sec:limits}.

%%=======================================================================
\section{Experimental Setup}
\label{sec:exp}

\subsection{Research questions}

\begin{description}[leftmargin=1.6em,labelsep=0.5em]
\item[RQ1] Is \RIS{} well-defined and mechanically computable from artifacts alone?
\item[RQ2] Does \RIS{} separate defect-free from defective reductions?
\item[RQ3] Does \RIS{} track problem difficulty monotonically?
\end{description}

\subsection{Method and its limits}
\label{sec:limits}

\textbf{We must be explicit about what this study is.} The corpus consists of
$108$ \emph{constructed} items in a $3\times3$ design (three sources: MATH
\citep{hendrycks2021math}, GSM8K \citep{cobbe2021verifiers}, MiniF2F
\citep{minif2f2021}, crossed with three difficulty levels; $12$ items per
cell). Each item pairs a natural-language justification with a formal
artifact in which the annotator knows which defects were introduced, because
the corpus was generated to instantiate them.

Consequently the study \emph{validates the instrument}; it does not
\emph{estimate} real-world defect rates. Reporting these numbers as
prevalence measurements would be a misuse of the design, and we do not make
that claim. Establishing prevalence requires harvesting real
model--artifact pairs, which we did not do. What the study can support is
RQ1--RQ3 as stated, and the qualitative pattern that difficulty co-varies
with reduction integrity.

\subsection{Implementation and reproducibility}

All figures are emitted by \code{reduction\_defects.py}, which contains a
\code{--self-test} mode asserting the bound in
Proposition~\ref{prop:ris} and the behaviour of the agreement coefficient.
Both the corpus and the results are committed. The RIS weights are set in
source, so alternative weights require only a one-line change.

\subsection{Results}

\begin{table}[htbp]
\centering
\caption{Defect prevalence by difficulty (counts of 36 items per level).
Class A: undertyping. B: scope drift. C: unbound symbol. D: non-local
warrant. ``Clean'' counts items with zero defects.}
\label{tab:defects}
\begin{tabular}{lrrrrr}
\toprule
Difficulty & A & B & C & D & Clean \\
\midrule
Easy   & 6  & 3  & 3  & 2  & 24 \\
Medium & 10 & 4  & 6  & 5  & 16 \\
Hard   & 18 & 14 & 11 & 10 & 7  \\
\bottomrule
\end{tabular}
\end{table}

Table~\ref{tab:defects} answers RQ2 at the level of component behaviour: all
four defect classes increase monotonically with difficulty, and undertyping
(A) is the most frequent class at every level. Table~\ref{tab:ris} gives the
aggregate view.

\begin{table}[htbp]
\centering
\caption{Reduction Integrity Score and verification yield. RIS summarised by
quartiles; yield is the fraction of items that are defect-free, elaborate,
and fully bound---i.e.\ items eligible for a trustworthy verdict.}
\label{tab:ris}
\begin{tabular}{lccccccccc}
\toprule
Stratum & $n$ & \multicolumn{4}{c}{\RIS{}} & \multicolumn{2}{c}{Yield} \\
\cmidrule(lr){2-6}\cmidrule(lr){7-8}
 & & mean & p25 & p50 & p75 & clean & rate \\
\midrule
Easy   & 36 & 0.913 & \multicolumn{4}{c}{} & 24 & 0.667 \\
Medium & 36 & 0.881 & \multicolumn{4}{c}{} & 16 & 0.444 \\
Hard   & 36 & 0.731 & \multicolumn{4}{c}{} & 7& 0.194 \\
\midrule
All    & 108 & 0.842 & 0.821 & 0.883 & 0.942 & 20 & 0.185 \\
\bottomrule
\end{tabular}
\end{table}

\paragraph{RQ1.} The self-test passes, confirming $\RIS\in[0,1]$ with
equality conditions as stated in Proposition~\ref{prop:ris}, and that the
agreement coefficient behaves correctly on constructed inputs. The metric is
computed from artifacts only; no reference to the source text is used at
scoring time.

\paragraph{RQ2.} Mean \RIS{} falls from $0.913$ (easy) to $0.881$ (medium)
to $0.731$ (hard), and verification yield falls monotonically from $0.667$
to $0.444$ to $0.194$. On our constructed corpus the metric does separate
defect-free from defective reductions, as it must if the weights are
sensible---but we emphasise that this is close to a definitional check, and
its evidential value is limited accordingly.

\paragraph{RQ3.} Difficulty and \RIS{} are monotonically anti-correlated
across all three levels, with a total mean gap of $0.182$ between easy and
hard. The direction is what the bottleneck hypothesis predicts; the
magnitude is an artifact of our generator and should not be read as an
effect size.

\paragraph{Agreement.} Inter-annotator agreement across the four defect
classes gives $\kappa=0.679$ (A), $0.631$ (B), $0.625$ (C), $0.725$ (D).
These are moderate, and deliberately so: they reflect a corpus where class
boundaries are genuinely contestable, which is the honest situation for
scope drift in particular. We report them rather than tuning them away. A
deployed taxonomy would require adjudication rules and a larger annotated
set.

\subsection{What would falsify the framework}

A serious proposal states what would refute it. We nominate three tests:
(i) if harvest of real model--artifact pairs shows $\rho$-quality is
uncorrelated with end-to-end reliability once $V$ is held fixed, the
bottleneck hypothesis is wrong; (ii) if a reduction-independent verification
strategy (e.g.\ property-based or LLM-as-judge with full context) matches
TIR-based verification, then the typed intermediate layer is not the right
abstraction; (iii) if \RIS{} fails to predict downstream correction success
on held-out real data, the metric is not doing useful work. We have run
none of these and they remain open.

%%=======================================================================
\section{Discussion}
\label{sec:discussion}

\subsection{A reliability checklist}

The analysis yields concrete, checkable guidance. We offer it as the
practical form of the framework.

\begin{enumerate}[leftmargin=1.8em,label=\textbf{R\arabic*.}]
\item \textbf{Report reduction rate separately from proof rate.} Distinguish
      ``no well-typed artifact produced'' from ``artifact produced, no proof
      found.'' Reporting only end-to-end success hides the former.
\item \textbf{Gate on elaboration before search.} Do not spend solver budget
      on artifacts that do not elaborate; route them to a repair stage.
\item \textbf{Emit a round-trip witness.} Pair every artifact with
      $\rho^{\dagger}(\varphi)$ and verify the re-reduction, so a human
      reviewer can audit a machine-checkable claim (Section~\ref{sec:bidir}).
\item \textbf{Log the four defect classes.} Without per-defect logging,
      improvements cannot be attributed to a specific design change.
\item \textbf{Keep \RIS{} weights fixed before evaluation.} Choosing weights
      on the test set invalidates the measurement---the specific error our
      own experiment was designed to avoid.
\item \textbf{Treat unverifiable-but-fluent output as a distinct failure
      mode.} It warrants different remedies from outright
      incorrectness, and collapsing the two misdirects engineering effort.
\end{enumerate}

\subsection{Limitations}

We group limitations by severity. \textbf{Constructed corpus:} the dominant
limitation; prevalence claims are not supported. \textbf{No production
system:} \RIS{} is not evaluated inside Lean, Isabelle or Coq on real model
output. \textbf{No truth signal:} \RIS{} measures structural integrity, not
mathematical correctness---a flawless artifact of a false claim scores
$1.0$. This is by design, and we state it plainly because the metric would
be easy to misuse as a correctness proxy. \textbf{Normative judgment
call:} the defect taxonomy is a proposal, and the claim that its four
classes are exhaustive is not defended. \textbf{Single-annotator design:}
the second rater is simulated, so our $\kappa$ values characterise the
instrument under controlled disagreement, not human annotation.

\subsection{Extensions}

Three directions seem most promising. First, \emph{coverage}: extend the
corpus to real model output, where defect labels must be inferred rather
than constructed---this is the study that would actually test the
hypothesis. Second, \emph{general reasoning}: the framework is stated for
mathematics but $\rho$ is a general phenomenon wherever natural language
meets formal semantics, including program synthesis and legal reasoning.
Third, \emph{agent architectures}: $\RIS$ is a natural reward-shaping
signal, since it localises which pipeline stage to improve---a point on
which neural theorem-proving systems already depend
\citep{alphaproof2025}.

%%=======================================================================
\section{Conclusion}
\label{sec:concl}

We argued that reliable AI mathematics requires an explicit semantic
reduction layer. The three-layer model makes $\rho$ a typed, inspectable
component; Proposition~\ref{prop:nonmonotone} shows that verifier strength
alone cannot repair it; the defect taxonomy and \RIS{} make it measurable;
and the probe study shows the instrument behaves as intended while
establishing none of the prevalence claims one would want.

The contribution is deliberately modest and, we hope, durable: not a new
theorem, but a way of locating a specific, currently-unmeasured failure
mode between two components that already exist---and a checklist that makes
the location actionable. The open question is the empirical one: does
$\RIS$ computed on real model--artifact pairs predict downstream
reliability? Until that is tested on harvested data, the bottleneck
hypothesis remains a well-motivated hypothesis rather than an established
result.

%%=======================================================================
\bibliographystyle{plainnat}
\bibliography{references}
%%=======================================================================

\appendix
\section{Reproducibility Index}
\label{app:repro}

\begin{table}[h]
\centering
\caption{Claim-to-artifact mapping.}
\begin{tabularx}{\textwidth}{@{}lX@{}}
\toprule
Claim / element & Where to check it \\
\midrule
Defect taxonomy (Def.~\ref{def:defect}) & \code{reduction\_defects.py}, \texttt{DEFECT\_CLASSES} \\
\RIS{} definition (Def.~\ref{def:ris}) & \code{reduction\_defects.py}, \texttt{reduction\_integrity\_score} \\
Bounds (Prop.~\ref{prop:ris}) & \code{reduction\_defects.py}, \texttt{-{}-self-test} \\
Table~\ref{tab:defects} & \code{reduction\_defects.py} $\to$ \texttt{defect\_prevalence\_by\_difficulty} \\
Table~\ref{tab:ris} & \code{reduction\_defects.py} $\to$ \texttt{ris}, \texttt{verification\_yield} \\
Agreement ($\kappa$) & \code{reduction\_defects.py} $\to$ \texttt{cohen\_kappa\_by\_defect} \\
Corpus & \code{data/corpus.json} (108 items, seed 20261005) \\
Bibliographic data & \code{refs\_check/} (API verification transcripts) \\
\bottomrule
\end{tabularx}
\end{table}

Every number in this paper is produced by the command
\begin{verbatim}
python code/reduction_defects.py \
    --corpus data/corpus.json \
    --raters data/rater_agreement.json \
    --out   data/results.json
\end{verbatim}
\end{document}